\chapter*{Annexe technique : code}
\addcontentsline{toc}{chapter}{Annexe technique : code}

\subsection{Le code suivant présente le script Python utilisé pour nettoyer les textes extraits des fichiers EPUB après leur conversion au format \textsl{.txt}}.

\begin{minted}[
    fontsize=\scriptsize,
    breaklines=true,
    breakanywhere=true,
    linenos=true,
    numbersep=5pt,
    frame=lines,
    framesep=2mm,
    baselinestretch=1,
    autogobble
]{python}
from pathlib import Path
import re


def nettoyage_structure_epub(texte):
    texte_propre = texte

    # Nettoyage initial des espaces insécables
    texte_propre = re.sub(r'[\xa0\u202f]', ' ', texte_propre)

    # Suppression des éléments placés entre crochets
    texte_propre = re.sub(r'\[.*?\]', '', texte_propre)

    # Conservation uniquement des caractères utiles pour l'analyse
    texte_propre = re.sub(
        r'[^a-zA-Z0-9àâäéèêëîïôöùûüçÀÂÄÉÈÊËÎÏÔÖÙÛÜÇœŒ'
        r'\s.,;:!?\-()\[\]]',
        '',
        texte_propre
    )

    # Suppression des chiffres collés aux mots
    texte_propre = re.sub(
        r'([a-zA-ZàâäéèêëîïôöùûüçÀÂÄÉÈÊËÎÏÔÖÙÛÜÇœŒ])\d+',
        r'\1',
        texte_propre
    )

    # Nettoyage des tirets et des guillemets
    texte_propre = re.sub(r'-{2,}', '-', texte_propre)
    texte_propre = re.sub(r'["«»“”]', '', texte_propre)

    # Suppression des numéros de page isolés
    texte_propre = re.sub(r'(?m)^\s*\d+\s*$', '', texte_propre)

    # Suppression des titres de chapitre en chiffres romains
    texte_propre = re.sub(r"(?im)^\s*[IVXLCDM]+\s*$", "", texte_propre)

    # Suppression de certains éléments propres aux fichiers EPUB
    texte_propre = re.sub(r"(?im)^\s*Liste des titres\s*$", "", texte_propre)
    texte_propre = re.sub(r"(?im)^\s*Table des matières du titre\s*$", "", texte_propre)

    # Suppression des lignes entièrement écrites en majuscules
    texte_propre = re.sub(
        r"(?m)^\s*[A-ZÀÂÄÉÈÊËÎÏÔÖÙÛÜÇ'’\s]+\s*$",
        "", texte_propre
    )

    # Transformation du texte en un seul bloc continu
    texte_propre = re.sub(r'\s+', ' ', texte_propre)

    # Segmentation du texte en phrases
    pattern_segmentation = r"(?<!\bM)([.?!])\s+(?=[A-ZÀÂÄÉÈÊËÎÏÔÖÙÛÜÇ\-])"
    texte_propre = re.sub(pattern_segmentation, r"\1\n", texte_propre)

    # Remise à la ligne des dialogues
    texte_propre = re.sub(r'\s+(?=-)', r'\n', texte_propre)

    return texte_propre.strip()


# Définition des dossiers source et destination
dossier_source = Path("nom_dossier")
dossier_destination = Path("Text_nettoyé/nom_dossier")

# Création du dossier de destination s'il n'existe pas
dossier_destination.mkdir(parents=True, exist_ok=True)

# Compteur permettant de suivre le nombre de fichiers traités
fichiers_traites = 0

print(f"Début du traitement des fichiers dans : {dossier_source}...\n")

# Boucle sur tous les fichiers .txt du dossier source
for fichier_source in dossier_source.glob("*.txt"):

    print(f"-> Nettoyage de : {fichier_source.name}")

    try:
        # Lecture du fichier brut
        contenu_brut = fichier_source.read_text(encoding="utf-8")

        # Nettoyage du texte
        contenu_propre = nettoyage_structure_epub(contenu_brut)

        # Création du nom du fichier nettoyé
        nom_fichier_propre = f"{fichier_source.stem}_clean.txt"
        fichier_destination = dossier_destination / nom_fichier_propre

        # Sauvegarde du fichier nettoyé
        fichier_destination.write_text(contenu_propre, encoding="utf-8")

        fichiers_traites += 1

    except Exception as e:
        print(f"   [ERREUR] Impossible de traiter {fichier_source.name} : {e}")

print(
    f"\nOpération terminée ! {fichiers_traites} fichiers ont été nettoyés "
    f"et placés dans le dossier '{dossier_destination}'."
)
\end{minted}